\documentclass[10pt,letterpaper]{article}
\usepackage{cornell-notes}
\usepackage{mathtools}

\title{Chapter 5: Evaluation of Functions}
\author{}
\date{}

\setDocTitle{Chapter 5: Evaluation of Functions}
\setDocSubtitle{Numerical Methods}
\setDocVersion{v1.0}
\setDocStatus{Study Companion}
\setDocOwner{Numerical Methods Cornell Notes}
\setDocAuthor{}
\setDocDate{\today}
\setDocSummary{Cornell-format study and review notes for Chapter 5: Evaluation of Functions.}

\setCornellCollection{Numerical Methods Cornell Notes}
\setCornellUnitType{Chapter}
\setCornellUnitNumber{5}
\setCornellUnitTitle{Evaluation of Functions}
\setCornellCompanionLabel{Study and review companion}

\hypersetup{pdftitle={Chapter 5: Evaluation of Functions},pdfsubject={Cornell notes},pdfauthor={}}

\begin{document}
\maketitle

\section*{Chapter Overview}
This chapter organizes the mathematical and computational ideas behind evaluation of functions. The notes preserve the numbered section sequence while emphasizing method selection, explicit assumptions, numerical reliability, cost, and verification.

\section*{Learning Objectives}
Explain the governing problem class; compare Polynomials and Rational Functions, Evaluation of Continued Fractions, Series and Their Convergence, and Recurrence Relations and Clenshaw's Recurrence Formula; design defensible stopping and verification checks; distinguish problem conditioning from algorithmic stability.

\section*{Prerequisites}
Series, polynomial algebra, recurrence relations, and floating-point arithmetic.

\section*{Operational Workflow}
Partition the domain; choose a stable representation for each region; scale arguments; evaluate with nested or recurrent forms; stop from an error estimate rather than equality of iterates.

\subsection*{Formula and notation anchor}
\[
p(x)=a_0+x\bigl(a_1+x(a_2+\cdots+x a_n)\bigr)
\]
Horner's rule uses $n$ multiplications and $n$ additions and usually avoids the extra rounding and cost of explicit powers.

\begin{warnbox}{Numerical warning}
Catastrophic cancellation, unstable recurrence direction, slow series, overflow in intermediate values, branch inconsistencies, and unsafe derivative step sizes.
\end{warnbox}

\section{5.0 Introduction}
\begin{cornell}
\noterow{What is the central idea?}{Function evaluation is an algorithm-design problem: algebraically equivalent formulas can differ greatly in cost, stability, and valid range.}
\noterow{How should it be used?}{For Introduction, begin from explicit inputs, outputs, preconditions, and representation choices.  Partition the domain; choose a stable representation for each region; scale arguments; evaluate with nested or recurrent forms; stop from an error estimate rather than equality of iterates.}
\noterow{Cost and reliability}{Analyze arithmetic work, storage, data movement, and convergence separately.  Relevant failure pressures include catastrophic cancellation, unstable recurrence direction, slow series, overflow in intermediate values, branch inconsistencies, and unsafe derivative step sizes.  Use scaled tolerances and report both success criteria and diagnostic evidence.}
\noterow{Implementation checklist}{\begin{itemize}
\item Validate dimensions, domains, ordering, and structural assumptions.
\item Guard small divisors, invalid states, overflow, underflow, and nonfinite values.
\item Make mutation, workspace, indexing, and termination behavior explicit.
\item Test a simple exact case, a difficult valid case, a boundary case, and an expected failure.
\end{itemize}}
\end{cornell}
\begin{summarybox}{Section summary}
Function evaluation is an algorithm-design problem: algebraically equivalent formulas can differ greatly in cost, stability, and valid range.  Select this technique only when its assumptions match the data, and accept a result only after an independent residual, error estimate, invariant, or refinement check.
\end{summarybox}

\section{5.1 Polynomials and Rational Functions}
\begin{cornell}
\noterow{What is the central idea?}{Nested evaluation minimizes polynomial operations, while rational functions require denominator checks and careful scaling.}
\noterow{How should it be used?}{For Polynomials and Rational Functions, begin from explicit inputs, outputs, preconditions, and representation choices.  Partition the domain; choose a stable representation for each region; scale arguments; evaluate with nested or recurrent forms; stop from an error estimate rather than equality of iterates.}
\noterow{Quantitative lens}{Identify the governing equation, recurrence, probability law, or geometric predicate for Polynomials and Rational Functions.  Define every variable and scale; then derive a residual, error estimate, or invariant that can be checked independently.}
\noterow{Cost and reliability}{Analyze arithmetic work, storage, data movement, and convergence separately.  Relevant failure pressures include catastrophic cancellation, unstable recurrence direction, slow series, overflow in intermediate values, branch inconsistencies, and unsafe derivative step sizes.  Use scaled tolerances and report both success criteria and diagnostic evidence.}
\noterow{Implementation checklist}{\begin{itemize}
\item Validate dimensions, domains, ordering, and structural assumptions.
\item Guard small divisors, invalid states, overflow, underflow, and nonfinite values.
\item Make mutation, workspace, indexing, and termination behavior explicit.
\item Test a simple exact case, a difficult valid case, a boundary case, and an expected failure.
\end{itemize}}
\end{cornell}
\begin{summarybox}{Section summary}
Nested evaluation minimizes polynomial operations, while rational functions require denominator checks and careful scaling.  Select this technique only when its assumptions match the data, and accept a result only after an independent residual, error estimate, invariant, or refinement check.
\end{summarybox}

\section{5.2 Evaluation of Continued Fractions}
\begin{cornell}
\noterow{What is the central idea?}{Continued fractions are evaluated through stable recurrences with safeguards against zero or overflowing intermediate denominators.}
\noterow{How should it be used?}{For Evaluation of Continued Fractions, begin from explicit inputs, outputs, preconditions, and representation choices.  Partition the domain; choose a stable representation for each region; scale arguments; evaluate with nested or recurrent forms; stop from an error estimate rather than equality of iterates.}
\noterow{Quantitative lens}{Identify the governing equation, recurrence, probability law, or geometric predicate for Evaluation of Continued Fractions.  Define every variable and scale; then derive a residual, error estimate, or invariant that can be checked independently.}
\noterow{Cost and reliability}{Analyze arithmetic work, storage, data movement, and convergence separately.  Relevant failure pressures include catastrophic cancellation, unstable recurrence direction, slow series, overflow in intermediate values, branch inconsistencies, and unsafe derivative step sizes.  Use scaled tolerances and report both success criteria and diagnostic evidence.}
\noterow{Implementation checklist}{\begin{itemize}
\item Validate dimensions, domains, ordering, and structural assumptions.
\item Guard small divisors, invalid states, overflow, underflow, and nonfinite values.
\item Make mutation, workspace, indexing, and termination behavior explicit.
\item Test a simple exact case, a difficult valid case, a boundary case, and an expected failure.
\end{itemize}}
\end{cornell}
\begin{summarybox}{Section summary}
Continued fractions are evaluated through stable recurrences with safeguards against zero or overflowing intermediate denominators.  Select this technique only when its assumptions match the data, and accept a result only after an independent residual, error estimate, invariant, or refinement check.
\end{summarybox}

\section{5.3 Series and Their Convergence}
\begin{cornell}
\noterow{What is the central idea?}{A series evaluator needs a convergence model, a defensible stopping test, and a strategy for slow convergence or cancellation.}
\noterow{How should it be used?}{For Series and Their Convergence, begin from explicit inputs, outputs, preconditions, and representation choices.  Partition the domain; choose a stable representation for each region; scale arguments; evaluate with nested or recurrent forms; stop from an error estimate rather than equality of iterates.}
\noterow{Quantitative lens}{Identify the governing equation, recurrence, probability law, or geometric predicate for Series and Their Convergence.  Define every variable and scale; then derive a residual, error estimate, or invariant that can be checked independently.}
\noterow{Cost and reliability}{Analyze arithmetic work, storage, data movement, and convergence separately.  Relevant failure pressures include catastrophic cancellation, unstable recurrence direction, slow series, overflow in intermediate values, branch inconsistencies, and unsafe derivative step sizes.  Use scaled tolerances and report both success criteria and diagnostic evidence.}
\noterow{Implementation checklist}{\begin{itemize}
\item Validate dimensions, domains, ordering, and structural assumptions.
\item Guard small divisors, invalid states, overflow, underflow, and nonfinite values.
\item Make mutation, workspace, indexing, and termination behavior explicit.
\item Test a simple exact case, a difficult valid case, a boundary case, and an expected failure.
\end{itemize}}
\end{cornell}
\begin{summarybox}{Section summary}
A series evaluator needs a convergence model, a defensible stopping test, and a strategy for slow convergence or cancellation.  Select this technique only when its assumptions match the data, and accept a result only after an independent residual, error estimate, invariant, or refinement check.
\end{summarybox}

\section{5.4 Recurrence Relations and Clenshaw's Recurrence Formula}
\begin{cornell}
\noterow{What is the central idea?}{Recurrences can evaluate expansions efficiently, but the stable direction depends on which solution of the recurrence dominates rounding errors.}
\noterow{How should it be used?}{For Recurrence Relations and Clenshaw's Recurrence Formula, begin from explicit inputs, outputs, preconditions, and representation choices.  Partition the domain; choose a stable representation for each region; scale arguments; evaluate with nested or recurrent forms; stop from an error estimate rather than equality of iterates.}
\noterow{Quantitative lens}{Identify the governing equation, recurrence, probability law, or geometric predicate for Recurrence Relations and Clenshaw's Recurrence Formula.  Define every variable and scale; then derive a residual, error estimate, or invariant that can be checked independently.}
\noterow{Cost and reliability}{Analyze arithmetic work, storage, data movement, and convergence separately.  Relevant failure pressures include catastrophic cancellation, unstable recurrence direction, slow series, overflow in intermediate values, branch inconsistencies, and unsafe derivative step sizes.  Use scaled tolerances and report both success criteria and diagnostic evidence.}
\noterow{Implementation checklist}{\begin{itemize}
\item Validate dimensions, domains, ordering, and structural assumptions.
\item Guard small divisors, invalid states, overflow, underflow, and nonfinite values.
\item Make mutation, workspace, indexing, and termination behavior explicit.
\item Test a simple exact case, a difficult valid case, a boundary case, and an expected failure.
\end{itemize}}
\end{cornell}
\begin{summarybox}{Section summary}
Recurrences can evaluate expansions efficiently, but the stable direction depends on which solution of the recurrence dominates rounding errors.  Select this technique only when its assumptions match the data, and accept a result only after an independent residual, error estimate, invariant, or refinement check.
\end{summarybox}

\section{5.5 Complex Arithmetic}
\begin{cornell}
\noterow{What is the central idea?}{Complex operations require branch conventions and scaled formulas, particularly for division, magnitude, roots, and transcendental functions.}
\noterow{How should it be used?}{For Complex Arithmetic, begin from explicit inputs, outputs, preconditions, and representation choices.  Partition the domain; choose a stable representation for each region; scale arguments; evaluate with nested or recurrent forms; stop from an error estimate rather than equality of iterates.}
\noterow{Quantitative lens}{Identify the governing equation, recurrence, probability law, or geometric predicate for Complex Arithmetic.  Define every variable and scale; then derive a residual, error estimate, or invariant that can be checked independently.}
\noterow{Cost and reliability}{Analyze arithmetic work, storage, data movement, and convergence separately.  Relevant failure pressures include catastrophic cancellation, unstable recurrence direction, slow series, overflow in intermediate values, branch inconsistencies, and unsafe derivative step sizes.  Use scaled tolerances and report both success criteria and diagnostic evidence.}
\noterow{Implementation checklist}{\begin{itemize}
\item Validate dimensions, domains, ordering, and structural assumptions.
\item Guard small divisors, invalid states, overflow, underflow, and nonfinite values.
\item Make mutation, workspace, indexing, and termination behavior explicit.
\item Test a simple exact case, a difficult valid case, a boundary case, and an expected failure.
\end{itemize}}
\end{cornell}
\begin{summarybox}{Section summary}
Complex operations require branch conventions and scaled formulas, particularly for division, magnitude, roots, and transcendental functions.  Select this technique only when its assumptions match the data, and accept a result only after an independent residual, error estimate, invariant, or refinement check.
\end{summarybox}

\section{5.6 Quadratic and Cubic Equations}
\begin{cornell}
\noterow{What is the central idea?}{Root formulas should be rearranged to avoid subtracting nearly equal quantities and should recover companion roots from algebraic invariants.}
\noterow{How should it be used?}{For Quadratic and Cubic Equations, begin from explicit inputs, outputs, preconditions, and representation choices.  Partition the domain; choose a stable representation for each region; scale arguments; evaluate with nested or recurrent forms; stop from an error estimate rather than equality of iterates.}
\noterow{Quantitative lens}{Identify the governing equation, recurrence, probability law, or geometric predicate for Quadratic and Cubic Equations.  Define every variable and scale; then derive a residual, error estimate, or invariant that can be checked independently.}
\noterow{Cost and reliability}{Analyze arithmetic work, storage, data movement, and convergence separately.  Relevant failure pressures include catastrophic cancellation, unstable recurrence direction, slow series, overflow in intermediate values, branch inconsistencies, and unsafe derivative step sizes.  Use scaled tolerances and report both success criteria and diagnostic evidence.}
\noterow{Implementation checklist}{\begin{itemize}
\item Validate dimensions, domains, ordering, and structural assumptions.
\item Guard small divisors, invalid states, overflow, underflow, and nonfinite values.
\item Make mutation, workspace, indexing, and termination behavior explicit.
\item Test a simple exact case, a difficult valid case, a boundary case, and an expected failure.
\end{itemize}}
\end{cornell}
\begin{summarybox}{Section summary}
Root formulas should be rearranged to avoid subtracting nearly equal quantities and should recover companion roots from algebraic invariants.  Select this technique only when its assumptions match the data, and accept a result only after an independent residual, error estimate, invariant, or refinement check.
\end{summarybox}

\section{5.7 Numerical Derivatives}
\begin{cornell}
\noterow{What is the central idea?}{Finite differences balance truncation error against cancellation; the best step is neither arbitrarily large nor arbitrarily small.}
\noterow{How should it be used?}{For Numerical Derivatives, begin from explicit inputs, outputs, preconditions, and representation choices.  Partition the domain; choose a stable representation for each region; scale arguments; evaluate with nested or recurrent forms; stop from an error estimate rather than equality of iterates.}
\noterow{Quantitative anchor}{\[
f'(x)\approx\frac{f(x+h)-f(x-h)}{2h}=f'(x)+O(h^2)
\]
State the dimensions, scaling, and validity assumptions before using this relation.  Check the resulting residual or invariant rather than trusting an iteration count alone.}
\noterow{Cost and reliability}{Analyze arithmetic work, storage, data movement, and convergence separately.  Relevant failure pressures include catastrophic cancellation, unstable recurrence direction, slow series, overflow in intermediate values, branch inconsistencies, and unsafe derivative step sizes.  Use scaled tolerances and report both success criteria and diagnostic evidence.}
\noterow{Implementation checklist}{\begin{itemize}
\item Validate dimensions, domains, ordering, and structural assumptions.
\item Guard small divisors, invalid states, overflow, underflow, and nonfinite values.
\item Make mutation, workspace, indexing, and termination behavior explicit.
\item Test a simple exact case, a difficult valid case, a boundary case, and an expected failure.
\end{itemize}}
\end{cornell}
\begin{summarybox}{Section summary}
Finite differences balance truncation error against cancellation; the best step is neither arbitrarily large nor arbitrarily small.  Select this technique only when its assumptions match the data, and accept a result only after an independent residual, error estimate, invariant, or refinement check.
\end{summarybox}

\section{5.8 Chebyshev Approximation}
\begin{cornell}
\noterow{What is the central idea?}{Chebyshev bases control worst-case polynomial error and support efficient approximation on a scaled finite interval.}
\noterow{How should it be used?}{For Chebyshev Approximation, begin from explicit inputs, outputs, preconditions, and representation choices.  Partition the domain; choose a stable representation for each region; scale arguments; evaluate with nested or recurrent forms; stop from an error estimate rather than equality of iterates.}
\noterow{Quantitative anchor}{\[
T_k(\cos\theta)=\cos(k\theta)
\]
State the dimensions, scaling, and validity assumptions before using this relation.  Check the resulting residual or invariant rather than trusting an iteration count alone.}
\noterow{Cost and reliability}{Analyze arithmetic work, storage, data movement, and convergence separately.  Relevant failure pressures include catastrophic cancellation, unstable recurrence direction, slow series, overflow in intermediate values, branch inconsistencies, and unsafe derivative step sizes.  Use scaled tolerances and report both success criteria and diagnostic evidence.}
\noterow{Implementation checklist}{\begin{itemize}
\item Validate dimensions, domains, ordering, and structural assumptions.
\item Guard small divisors, invalid states, overflow, underflow, and nonfinite values.
\item Make mutation, workspace, indexing, and termination behavior explicit.
\item Test a simple exact case, a difficult valid case, a boundary case, and an expected failure.
\end{itemize}}
\end{cornell}
\begin{summarybox}{Section summary}
Chebyshev bases control worst-case polynomial error and support efficient approximation on a scaled finite interval.  Select this technique only when its assumptions match the data, and accept a result only after an independent residual, error estimate, invariant, or refinement check.
\end{summarybox}

\section{5.9 Derivatives or Integrals of a Chebyshev-Approximated Function}
\begin{cornell}
\noterow{What is the central idea?}{Coefficient recurrences differentiate or integrate a Chebyshev expansion while retaining a compact basis representation.}
\noterow{How should it be used?}{For Derivatives or Integrals of a Chebyshev-Approximated Function, begin from explicit inputs, outputs, preconditions, and representation choices.  Partition the domain; choose a stable representation for each region; scale arguments; evaluate with nested or recurrent forms; stop from an error estimate rather than equality of iterates.}
\noterow{Quantitative anchor}{\[
T_k(\cos\theta)=\cos(k\theta)
\]
State the dimensions, scaling, and validity assumptions before using this relation.  Check the resulting residual or invariant rather than trusting an iteration count alone.}
\noterow{Cost and reliability}{Analyze arithmetic work, storage, data movement, and convergence separately.  Relevant failure pressures include catastrophic cancellation, unstable recurrence direction, slow series, overflow in intermediate values, branch inconsistencies, and unsafe derivative step sizes.  Use scaled tolerances and report both success criteria and diagnostic evidence.}
\noterow{Implementation checklist}{\begin{itemize}
\item Validate dimensions, domains, ordering, and structural assumptions.
\item Guard small divisors, invalid states, overflow, underflow, and nonfinite values.
\item Make mutation, workspace, indexing, and termination behavior explicit.
\item Test a simple exact case, a difficult valid case, a boundary case, and an expected failure.
\end{itemize}}
\end{cornell}
\begin{summarybox}{Section summary}
Coefficient recurrences differentiate or integrate a Chebyshev expansion while retaining a compact basis representation.  Select this technique only when its assumptions match the data, and accept a result only after an independent residual, error estimate, invariant, or refinement check.
\end{summarybox}

\section{5.10 Polynomial Approximation from Chebyshev Coefficients}
\begin{cornell}
\noterow{What is the central idea?}{A truncated Chebyshev series can be converted or evaluated as a polynomial, with truncation chosen from coefficient decay and target error.}
\noterow{How should it be used?}{For Polynomial Approximation from Chebyshev Coefficients, begin from explicit inputs, outputs, preconditions, and representation choices.  Partition the domain; choose a stable representation for each region; scale arguments; evaluate with nested or recurrent forms; stop from an error estimate rather than equality of iterates.}
\noterow{Quantitative anchor}{\[
T_k(\cos\theta)=\cos(k\theta)
\]
State the dimensions, scaling, and validity assumptions before using this relation.  Check the resulting residual or invariant rather than trusting an iteration count alone.}
\noterow{Cost and reliability}{Analyze arithmetic work, storage, data movement, and convergence separately.  Relevant failure pressures include catastrophic cancellation, unstable recurrence direction, slow series, overflow in intermediate values, branch inconsistencies, and unsafe derivative step sizes.  Use scaled tolerances and report both success criteria and diagnostic evidence.}
\noterow{Implementation checklist}{\begin{itemize}
\item Validate dimensions, domains, ordering, and structural assumptions.
\item Guard small divisors, invalid states, overflow, underflow, and nonfinite values.
\item Make mutation, workspace, indexing, and termination behavior explicit.
\item Test a simple exact case, a difficult valid case, a boundary case, and an expected failure.
\end{itemize}}
\end{cornell}
\begin{summarybox}{Section summary}
A truncated Chebyshev series can be converted or evaluated as a polynomial, with truncation chosen from coefficient decay and target error.  Select this technique only when its assumptions match the data, and accept a result only after an independent residual, error estimate, invariant, or refinement check.
\end{summarybox}

\section{5.11 Economization of Power Series}
\begin{cornell}
\noterow{What is the central idea?}{Economization replaces negligible high-order behavior by a lower-degree approximation that controls error over an interval rather than at one point.}
\noterow{How should it be used?}{For Economization of Power Series, begin from explicit inputs, outputs, preconditions, and representation choices.  Partition the domain; choose a stable representation for each region; scale arguments; evaluate with nested or recurrent forms; stop from an error estimate rather than equality of iterates.}
\noterow{Quantitative lens}{Identify the governing equation, recurrence, probability law, or geometric predicate for Economization of Power Series.  Define every variable and scale; then derive a residual, error estimate, or invariant that can be checked independently.}
\noterow{Cost and reliability}{Analyze arithmetic work, storage, data movement, and convergence separately.  Relevant failure pressures include catastrophic cancellation, unstable recurrence direction, slow series, overflow in intermediate values, branch inconsistencies, and unsafe derivative step sizes.  Use scaled tolerances and report both success criteria and diagnostic evidence.}
\noterow{Implementation checklist}{\begin{itemize}
\item Validate dimensions, domains, ordering, and structural assumptions.
\item Guard small divisors, invalid states, overflow, underflow, and nonfinite values.
\item Make mutation, workspace, indexing, and termination behavior explicit.
\item Test a simple exact case, a difficult valid case, a boundary case, and an expected failure.
\end{itemize}}
\end{cornell}
\begin{summarybox}{Section summary}
Economization replaces negligible high-order behavior by a lower-degree approximation that controls error over an interval rather than at one point.  Select this technique only when its assumptions match the data, and accept a result only after an independent residual, error estimate, invariant, or refinement check.
\end{summarybox}

\section{5.12 Pade Approximants}
\begin{cornell}
\noterow{What is the central idea?}{A Pade approximant matches a power series with a rational function, often extending usefulness beyond the convergence behavior of a polynomial truncation.}
\noterow{How should it be used?}{For Pade Approximants, begin from explicit inputs, outputs, preconditions, and representation choices.  Partition the domain; choose a stable representation for each region; scale arguments; evaluate with nested or recurrent forms; stop from an error estimate rather than equality of iterates.}
\noterow{Quantitative lens}{Identify the governing equation, recurrence, probability law, or geometric predicate for Pade Approximants.  Define every variable and scale; then derive a residual, error estimate, or invariant that can be checked independently.}
\noterow{Cost and reliability}{Analyze arithmetic work, storage, data movement, and convergence separately.  Relevant failure pressures include catastrophic cancellation, unstable recurrence direction, slow series, overflow in intermediate values, branch inconsistencies, and unsafe derivative step sizes.  Use scaled tolerances and report both success criteria and diagnostic evidence.}
\noterow{Implementation checklist}{\begin{itemize}
\item Validate dimensions, domains, ordering, and structural assumptions.
\item Guard small divisors, invalid states, overflow, underflow, and nonfinite values.
\item Make mutation, workspace, indexing, and termination behavior explicit.
\item Test a simple exact case, a difficult valid case, a boundary case, and an expected failure.
\end{itemize}}
\end{cornell}
\begin{summarybox}{Section summary}
A Pade approximant matches a power series with a rational function, often extending usefulness beyond the convergence behavior of a polynomial truncation.  Select this technique only when its assumptions match the data, and accept a result only after an independent residual, error estimate, invariant, or refinement check.
\end{summarybox}

\section{5.13 Rational Chebyshev Approximation}
\begin{cornell}
\noterow{What is the central idea?}{Rational Chebyshev approximations combine interval-aware polynomial bases with a denominator model suited to sharp or near-singular behavior.}
\noterow{How should it be used?}{For Rational Chebyshev Approximation, begin from explicit inputs, outputs, preconditions, and representation choices.  Partition the domain; choose a stable representation for each region; scale arguments; evaluate with nested or recurrent forms; stop from an error estimate rather than equality of iterates.}
\noterow{Quantitative anchor}{\[
T_k(\cos\theta)=\cos(k\theta)
\]
State the dimensions, scaling, and validity assumptions before using this relation.  Check the resulting residual or invariant rather than trusting an iteration count alone.}
\noterow{Cost and reliability}{Analyze arithmetic work, storage, data movement, and convergence separately.  Relevant failure pressures include catastrophic cancellation, unstable recurrence direction, slow series, overflow in intermediate values, branch inconsistencies, and unsafe derivative step sizes.  Use scaled tolerances and report both success criteria and diagnostic evidence.}
\noterow{Implementation checklist}{\begin{itemize}
\item Validate dimensions, domains, ordering, and structural assumptions.
\item Guard small divisors, invalid states, overflow, underflow, and nonfinite values.
\item Make mutation, workspace, indexing, and termination behavior explicit.
\item Test a simple exact case, a difficult valid case, a boundary case, and an expected failure.
\end{itemize}}
\end{cornell}
\begin{summarybox}{Section summary}
Rational Chebyshev approximations combine interval-aware polynomial bases with a denominator model suited to sharp or near-singular behavior.  Select this technique only when its assumptions match the data, and accept a result only after an independent residual, error estimate, invariant, or refinement check.
\end{summarybox}

\section{5.14 Evaluation of Functions by Path Integration}
\begin{cornell}
\noterow{What is the central idea?}{When derivatives satisfy an auxiliary differential equation, continuation along a safe path can evaluate a function outside an easy starting region.}
\noterow{How should it be used?}{For Evaluation of Functions by Path Integration, begin from explicit inputs, outputs, preconditions, and representation choices.  Partition the domain; choose a stable representation for each region; scale arguments; evaluate with nested or recurrent forms; stop from an error estimate rather than equality of iterates.}
\noterow{Quantitative lens}{Identify the governing equation, recurrence, probability law, or geometric predicate for Evaluation of Functions by Path Integration.  Define every variable and scale; then derive a residual, error estimate, or invariant that can be checked independently.}
\noterow{Cost and reliability}{Analyze arithmetic work, storage, data movement, and convergence separately.  Relevant failure pressures include catastrophic cancellation, unstable recurrence direction, slow series, overflow in intermediate values, branch inconsistencies, and unsafe derivative step sizes.  Use scaled tolerances and report both success criteria and diagnostic evidence.}
\noterow{Implementation checklist}{\begin{itemize}
\item Validate dimensions, domains, ordering, and structural assumptions.
\item Guard small divisors, invalid states, overflow, underflow, and nonfinite values.
\item Make mutation, workspace, indexing, and termination behavior explicit.
\item Test a simple exact case, a difficult valid case, a boundary case, and an expected failure.
\end{itemize}}
\end{cornell}
\begin{summarybox}{Section summary}
When derivatives satisfy an auxiliary differential equation, continuation along a safe path can evaluate a function outside an easy starting region.  Select this technique only when its assumptions match the data, and accept a result only after an independent residual, error estimate, invariant, or refinement check.
\end{summarybox}

\section*{Method comparison}
\begin{center}
\small
\begin{tabularx}{\linewidth}{>{\bfseries\raggedright\arraybackslash}p{0.22\linewidth}XX}
\toprule
Method or lens & Best use & Main caution \\
\midrule
Power or Taylor series & Near an expansion point & Can converge slowly or leave its radius \\
Chebyshev series & Uniform interval approximation & Requires scaling to the canonical interval \\
Rational approximation & Wide range or pole-like behavior & Denominator safety is essential \\
\bottomrule
\end{tabularx}
\end{center}

\section*{Worked example}
\begin{exambox}{Independent worked example}
Evaluate $p(x)=1-3x+2x^2$ at $x=1.2$ by Horner's rule: $p=1+x(-3+2x)=1+1.2(-0.6)=0.28$.  Direct evaluation gives the same mathematical value.  Horner's rule reduces operations, though any near-cancellation in the polynomial's value still limits relative accuracy.
\end{exambox}

\section*{Chapter synthesis}
\begin{summarybox}{Chapter synthesis}
The chapter's methods should be viewed as a decision system rather than a list of routines.  Begin with the mathematical problem and its conditioning, exploit structure, select an algorithm with compatible stability and cost, and verify the computed answer with evidence that is independent of the internal iteration.  The recurring implementation discipline is: partition the domain; choose a stable representation for each region; scale arguments; evaluate with nested or recurrent forms; stop from an error estimate rather than equality of iterates.
\end{summarybox}

\subsection*{Common mistakes and numerical hazards}
\begin{warnbox}{Common mistakes and numerical hazards}
Catastrophic cancellation, unstable recurrence direction, slow series, overflow in intermediate values, branch inconsistencies, and unsafe derivative step sizes.  A second recurring mistake is reporting only a computed value: always retain tolerances, iteration counts, residuals, refinement behavior, and relevant structural checks.
\end{warnbox}

\subsection*{Retrieval practice}
\textbf{Questions}
\begin{enumerate}
\item What problem does Polynomials and Rational Functions address, and which assumption is easiest to violate?
\item What problem does Evaluation of Continued Fractions address, and which assumption is easiest to violate?
\item What problem does Series and Their Convergence address, and which assumption is easiest to violate?
\item What problem does Rational Chebyshev Approximation address, and which assumption is easiest to violate?
\item What problem does Evaluation of Functions by Path Integration address, and which assumption is easiest to violate?
\item How do conditioning, stability, and the final residual play different roles in this chapter?
\end{enumerate}

\textbf{Brief answers}
\begin{enumerate}
\item Nested evaluation minimizes polynomial operations, while rational functions require denominator checks and careful scaling. Recheck the section's domain, structure, scaling, and failure diagnostics.
\item Continued fractions are evaluated through stable recurrences with safeguards against zero or overflowing intermediate denominators. Recheck the section's domain, structure, scaling, and failure diagnostics.
\item A series evaluator needs a convergence model, a defensible stopping test, and a strategy for slow convergence or cancellation. Recheck the section's domain, structure, scaling, and failure diagnostics.
\item Rational Chebyshev approximations combine interval-aware polynomial bases with a denominator model suited to sharp or near-singular behavior. Recheck the section's domain, structure, scaling, and failure diagnostics.
\item When derivatives satisfy an auxiliary differential equation, continuation along a safe path can evaluate a function outside an easy starting region. Recheck the section's domain, structure, scaling, and failure diagnostics.
\item Conditioning describes sensitivity of the mathematical problem, stability describes error propagation by the algorithm, and the residual measures satisfaction of the computed equations; none is a universal substitute for the others.
\end{enumerate}

\subsection*{Mastery checklist}
\begin{itemize}
\item[$\square$] I can explain and select Polynomials and Rational Functions without relying on a memorized code listing.
\item[$\square$] I can explain and select Evaluation of Continued Fractions without relying on a memorized code listing.
\item[$\square$] I can explain and select Series and Their Convergence without relying on a memorized code listing.
\item[$\square$] I can explain and select Recurrence Relations and Clenshaw's Recurrence Formula without relying on a memorized code listing.
\item[$\square$] I can explain and select Complex Arithmetic without relying on a memorized code listing.
\item[$\square$] I can state a scaled stopping rule and an independent verification check.
\item[$\square$] I can identify at least one conditioning risk and one algorithmic stability risk.
\item[$\square$] I can estimate time, storage, and data-movement costs at the appropriate level.
\end{itemize}

\end{document}

